% id: recurrence
% title: Recurrence of planar Brownian motion
% date: 05 Mar. 2026

\documentclass{article}
\usepackage{amsmath, amssymb}

\begin{document}

I first saw the recurrence (in small balls) property of planar Brownian motion a couple years ago during 
the Advanced Probability course offered by Part III. The proof I saw was of course (I believe) the 
standard approach in which one uses the fact that the Brownian motion composed with any harmonic function 
remains to be a martingale. In particular, by using the logarithm, this fact (and the optional stopping theorem) 
allows us to easily compute the probability that 
the Brownian motion hits a small ball before hitting a larger one. Thus, by taking the radius of the 
larger ball to infinity, we finds that the Brownian motion hits the small ball with probability one. 

This problem came up again recently during the exercise session of the Random Geometry course offered 
here at EPFL and, in an attempt at avoiding martingale techniques, Colin Piernot and I (mostly Colin)
found a clever proof of this the sketch of which I would like to share here. 

\begin{theorem}
  Let \(B_t\) be a complex Brownian motion and \(z\) be a point in the complex plane. Then, for any 
  \(\epsilon > 0\), we have that
  \[\mathbb{P}(|B_t - z| < \epsilon \text{ i.o.}) = 1.\]
\end{theorem}
\begin{proof}
  By translation, we my assume without loss of generality that \(z = 0\) and the Brownian motions 
  starts at \(1\).

  Let us define the following continuous time random walk \(X_t\) on \(\mathbb{Z}\): 
  Defining \(\tau_0 = 0\)
  \[\tau_{n + 1} := \inf\{t > \tau_n : |B_t| \in 2^{\mathbb{Z}}\}\]
  we set \(X_t = \log_2 |B_{\tau_n}|\) whenever \(t \in [\tau_n, \tau_{n + 1})\). Namely, 
  \(X_t\) is the continuous time random walk which moves to \(n\) whenever the Brownian motion 
  hits the circle of radius \(2^n\). Thus, as hitting the \(\epsilon\)-ball centered at the origin 
  corresponds to \(X_t\) hitting \(n\) for any \(n \in \mathbb{Z}\) satisfying \(2^n < \epsilon\), 
  it follows that the conclusion holds if we can show that \(X_t\) is symmetric (and so, is transient). 

  Straightaway, by scaling invariance and the strong Markov property, it suffices to show that 
  \[\mathbb{P}(X_1 = 0) = \mathbb{P}(X_1 = 2) = \frac{1}{2}.\]
  Namely, we need to show that \(B_t\) hits the circle of radius \(1/2\) before the circle of 
  radius \(2\) with probability \(1/2\).
  However, this is clear as, the map 
  \[f : A \to A : z \mapsto \frac{1}{z}\]
  with \(A\) being the annulus \(\{z \in \mathbb{C} : 1/2 \le |z| \le 2\}\) is conformal. Thus, 
  as \(f\) maps the boundary \(\{|z| = 2\}\) to \(\{|z| = 1/2\}\) and \(f(B_t)\) is also a Brownian 
  motion starting at 1 (by the conformal invariance of Brownian motion), the conclusion holds 
  by symmetry.
\end{proof}

Interestingly (or maybe not so considering the topic of the course), this type of argument also 
appeared later in the Random Geometry course in the probabilistic proof of the Picards little theorem.
There, starting with the Brownian motion's image over a conformal map, one similarly constructs a 
random walk on \(\mathbb{Z}^2\) (albeit here the construction is much more complicated with a lot 
of nuances. In particular, here the random walk is associated with the winding number of the trace 
quotiented by an equivalence relation in order to obtain integrability). Then, by showing that this 
radom walk is transient, one is able to conclude the theorem.

\end{document}
